\section{A no-teleportation theorem}
\label{sec:no-teleportation}

Four spaces recur in the block:
\begin{enumerate}[label=(\alph*),leftmargin=2.2em]
\item the actual nonzero provenance incidence graph
\(\cG_X^{\rm inc}\) of TPC-77/78;
\item the common synthesis target \(K_X\) and quotient
\(K_X/V_{A,X}\) of TPC-79;
\item the scalar row space on which the Hermitian pair
\((H_X^{\rm hol},P_X^{\rm hol})\) of TPC-80 acts; and
\item the determinant-bin space \(\C^{\cN_X}\) containing \(a_X\).
\end{enumerate}
They can share indices or provenance labels while carrying different
operations.

\begin{definition}[Literal bridge]
\label{def:literal-bridge}
A literal bridge between two audited objects is a specified linear or
affine map, constructed from the complete physical packet, that:
\begin{enumerate}[label=(B\arabic*),leftmargin=2.7em]
\item preserves all native multiplicities, signs, phases, and
normalization factors;
\item has the stated source and target, rather than an observation
space with the arrows reversed;
\item maps actual nonzero support to the actual object being
estimated; and
\item proves the norm, quadratic-form, or aggregation inequality
needed for the next step uniformly in the original quantifier range.
\end{enumerate}
\end{definition}

\begin{theorem}[No teleportation]
\label{thm:no-teleportation}
Data consisting only of
\[
 \cG_X^{\rm inc},\qquad K_X/V_{A,X},\qquad
 (H_X^{\rm hol},P_X^{\rm hol}),
\]
and their support, quotient-gauge, or scalar-spectral invariants do
not determine a nontrivial cancellation bound for \(|Z_X|\) or a
positive lower bound for \(D_X\).  Any inference from the listed
structural data to such a conclusion requires, at minimum, a literal
bridge
to the pre-bin coefficient and an exact determinant map
\[
 \cJ_X:\C^{\cY_X}\longrightarrow\C^{\cN_X}.
\]
In particular, the conclusions of TPC-78, TPC-79, and TPC-80 cannot
substitute for one another or for a determinant-frequency estimate.
\end{theorem}

\begin{proof}
The listed structural data contain no map to
\(\C^{\cN_X}\).  Keeping them fixed therefore imposes no algebraic
restriction on a vector assigned to that separate space.  One may
assign the same structural triple to \(a=(1,\ldots,1)\) or, for even
dimension, to \(b=(1,-1,\ldots,1,-1)\).  The two determinant vectors
have the same support, entrywise absolute values, and energy, while
their zero modes are different.  Thus the structural triple does not
determine \(Z_X\).

Even a map into a raw atomic space does not determine \(D_X\): a
determinant fiber may send two atoms \((1,-1)\) to zero.  Hence a
positive raw norm does not give a positive post-bin norm without a
lower modulus for the actual \(\cJ_X\) on a verified range.  These
examples are formalized in \cref{sec:counterexamples}.  A literal
bridge and determinant map add exactly the missing relations.
\end{proof}

\begin{corollary}[Three prohibited identifications]
\label{cor:prohibited-identifications}
Without additional proofs, none of the following is valid:
\begin{enumerate}[label=(\roman*),leftmargin=2.2em]
\item treating a four-corner observation or trace as an anchor
synthesis vector in the direct target;
\item treating a block-valued carrier as a scalar Hermitian edge
matrix, or replacing the actual row graph by an ambient graph; or
\item treating raw atomic mass or pre-bin energy as determinant-bin
energy.
\end{enumerate}
\end{corollary}

\begin{proof}
Item (i) reverses or changes the target of a map, item (ii) changes
the operator algebra, and item (iii) discards the kernel of
\(\cJ_X\).  Each violates \cref{def:literal-bridge}.
\end{proof}

\begin{remark}[What the theorem does not say]
\Cref{thm:no-teleportation} is a typed nonidentifiability result.  It
does not assert that no useful bridge exists.  Constructing such a
bridge is precisely Route~I-A of \cref{sec:decision-tree}; the direct
determinant-energy Route~I-B does not require this completion bridge.
\end{remark}
